% =====================================================================
% methods.tex — DRAFTED at rb-050 (RB-050), 2026-05-31.
%
% Fourth and final bookend; companion of sections/abstract.tex (rb-047),
% sections/intro.tex (rb-049), and sections/limitations.tex (rb-048).
% Catalogues the rebuild's simulation infrastructure (Rebuild/model/,
% Rebuild/sims/, Rebuild/derivations/) paralleling the original paper's
% §3 methods.  Closes the structural arc:
%   abstract → intro → model → results → extensions → limitations →
%   methods → appendix.
%
% Strength check (mission §3, CLAIM_LEDGER.md): no claim strength is
% changed here; the section names the infrastructure that already
% backs every claim in the body sections.  Recovery-test sha256
% digests and per-sim output digests are cited verbatim from
% rebuilder_state.json.
% =====================================================================

\section{Methods: the rebuilt simulation infrastructure}
\label{sec:methods}

This section catalogues the workspace under \texttt{Rebuild/} that
backs every quantitative claim in the body sections.  It parallels
the inherited paper's \S3 in organisation --- reference implementation,
simulation protocol, reproducibility --- and restates each in the
rebuild's voice: every extension to the inherited model is gated by a
recovery test that pins the extended model to the inherited model in
the appropriate limit to floating-point identity (or, where
floating-point identity is structurally impossible, to a tolerance no
coarser than $10^{-9}$); every simulation under \texttt{Rebuild/sims/}
records a deterministic SHA-256 digest of its canonical JSON output,
so the simulation is the operational definition of the claim it
backs.  No quantitative claim in the manuscript exceeds the strength
licensed by its row of \texttt{CLAIM\_LEDGER.md}, and every such row
points at one or more of the artifacts catalogued below.

% ---------------------------------------------------------------------
\subsection{The validated reference implementation}
\label{sec:methods-model}

The rebuilt model lives in \texttt{Rebuild/model/core.py} and is
documented in \texttt{Rebuild/model/README.md}.  Its starting point is
not a rewrite but a copy: the validated P1--P4 optimiser from the
skeptical-review replication \texttt{Critique/replications/\linebreak C5--symmetric-recovery/}
was copied into \texttt{Rebuild/model/} and extended; the inherited
model's headline numbers (Section~\ref{sec:model-inherited}) therefore
appear in the rebuilt model at byte-exact backwards compatibility
under the inherited parameter defaults, by construction.  The
extensions add four orthogonal axes, each gated by a recovery test
under \texttt{Rebuild/model/tests/}:

\begin{enumerate}
  \item \textbf{A1 --- the decorrelation channel.}  The per-location
    independence premise of the inherited model
    (Equation~\eqref{eq:pnofa-indep}) is replaced by the exact
    equicorrelated-Gaussian orthant integral
    (Equation~\eqref{eq:pnofa-rho}), evaluated by one-factor
    Gauss--Hermite quadrature at $n_q = 64$ nodes (accurate to
    $\sim 10^{-16}$ vs.\ $n_q = 128$ across the operating envelope).
    The inherited independent case is recovered exactly at
    $\corr = 0$ (Equation~\eqref{eq:rho-zero-recovery}; Definition
    \ref{def:three-levers} promotes $\corr$ to a first-class lever).
    Backs Sections \ref{sec:model-rho-channel} and
    \ref{sec:model-upper-bound}; full derivation in
    Appendix~\ref{sec:appendix-deriv-a1}.

  \item \textbf{A3 --- the power-mean conservation family.}  The
    additive conservation rule $\benefit + \cost = 2$ is generalised
    to the power-mean family $M_p(\benefit, \cost) = 1$ with closed-
    form weights $\cost(\Rsens; p) = (2/(\Rsens^p + 1))^{1/p}$ and
    $\benefit = \Rsens \cost$
    (Equations~\eqref{eq:conservation-family} and
    \eqref{eq:beta-gamma-of-p}).  The inherited additive case is
    $p = 1$; the multiplicative alternative
    $\benefit \cdot \cost = 1$ is the limit $p = 0$.  The full
    Hardy--Littlewood--P\'olya monotonicity is established in the
    rebuild's voice in
    Appendix~\ref{sec:appendix-deriv-a3} via the closed-form
    KL-divergence identity~\eqref{eq:a3d-hlp-kl}.  Backs
    Section~\ref{sec:extensions-a3}.

  \item \textbf{A2 --- per-location benefit/cost ratios.}  The
    single global scalar $\Rsens$ is promoted to a per-location vector
    $\Rsens_i$ via the heterogeneous-$\dprime$ map
    (Equation~\eqref{eq:d-prime-hetero}).  The inherited single-$\Rsens$
    model is recovered byte-for-byte under uniform $\Rsens_i$ and the
    canonical homogeneous allocation.  Backs
    Section~\ref{sec:extensions-a2}.

  \item \textbf{A8 --- the $N$-dimensional allocation simplex.}  The
    inherited 1-D allocation space (one $\alphacued$ for the cued
    slot, uniform $(1 - \alphacued)/(\Nloc - 1)$ on each uncued slot)
    is generalised to the full $\Nloc$-dimensional simplex
    $\sum_i \alpha_i = 1$, scored via the $\corr$-aware grouped-criterion
    optimiser that composes cleanly with all three prior extensions.
    The inherited homogeneous case is recovered to $10^{-9}$ absolute,
    the residual being the ULP-level reconstruction error
    $(\Nloc - 1) \cdot ((1 - \valid)/(\Nloc - 1)) - (1 - \valid)$;
    every headline number this composition drives is reported to no
    better than four decimal places, so $10^{-9}$ is six orders of
    magnitude past sensitive.  Backs
    Section~\ref{sec:extensions-a8}.
\end{enumerate}

Each axis composes with the others, and the joint composition is
itself a recovery contract: with $\corr = 0$, $p = 1$, uniform
$\Rsens_i$, and the canonical homogeneous allocation, the rebuilt
optimiser reproduces every inherited number in the
$4{,}410$-cell sweep at floating-point identity.  The notation in
\texttt{Rebuild/model/core.py} matches the manuscript verbatim:
\texttt{d\_prime\_asym} and \texttt{d\_prime\_hetero} are the
sensitivity maps of Sections~\ref{sec:model-inherited} and
\ref{sec:extensions-a2}; \texttt{beta\_gamma} is the
$(\benefit, \cost)$ pair of Equation~\eqref{eq:beta-gamma-of-p};
\texttt{p\_no\_fa\_indep} and \texttt{p\_no\_fa\_rho} are the
independent and equicorrelated $\PnoFA$ of
Equations~\eqref{eq:pnofa-indep} and \eqref{eq:pnofa-rho};
\texttt{policies} and \texttt{er\_full\_policy} score the four nested
policies $\Rpone, \Rpfour, \Rptwo, \Rpthree$ at the inherited 1-D and
the A8 $\Nloc$-D allocation spaces respectively.

% ---------------------------------------------------------------------
\subsection{Recovery contracts}
\label{sec:methods-recovery}

Four recovery tests under \texttt{Rebuild/model/tests/} are run before
any simulation on the rebuilt model is trusted.  Each test asserts
byte-exact (or, where structurally impossible, $10^{-9}$-tolerance)
reproduction of an inherited-model identity that the rebuild's
extension must preserve in the appropriate limit; each test produces
a deterministic SHA-256 digest of its JSON-serialised pass/fail
report.  Table~\ref{tab:methods-recovery} lists the four contracts
verbatim.

\begin{table}[tb]
  \centering
  \small
  \begin{tabular}{@{}lllp{6.5cm}@{}}
    \toprule
    Contract & Axis & Limit & SHA-256 digest (canonical JSON) \\
    \midrule
    \texttt{test\_recovery.py}
      & A1 ($\corr$ channel)
      & $\corr \to 0$
      & \texttt{d3c62215cedbe2f797c484a8d76e8bef1d5ac42cdb3092ff099a34aea1b6f18f} \\
    \texttt{test\_conservation\_family.py}
      & A3 (power mean $p$)
      & $p \to 1$
      & \texttt{f4f57a89e5108db47447cbeaf2f15440f56cdfffb65f3d173dc4a0550121791e} \\
    \texttt{test\_heterogeneous\_r.py}
      & A2 (per-loc $\Rsens_i$)
      & uniform $\Rsens_i$
      & \texttt{0486921f8a4e253a1682ee0731e52c7de800bf6c3dbb1b5e8ff3200602454c7e} \\
    \texttt{test\_general\_policy.py}
      & A8 ($\Nloc$-D simplex)
      & homogeneous uncued
      & \texttt{883ea15af9fd069e04c05ff156d65f33a7d25278891092539c6441d2248c3d39} \\
    \bottomrule
  \end{tabular}
  \caption[Recovery contracts]{The four recovery contracts gating the
    rebuilt model's extension axes.  Each contract recovers the
    inherited model in the named limit; the SHA-256 digest is over
    the canonical (sort-keys, indent-$2$) JSON serialisation of the
    test's pass/fail report.  A simulation is run only after its
    contributing extensions are pinned by their corresponding
    contracts.}
  \label{tab:methods-recovery}
\end{table}

These four digests appear verbatim in
\texttt{Rebuild/rebuilder\_state.json} as the fields
\texttt{recovery\_test\_digest},
\texttt{conservation\_family\_test\_digest},
\texttt{heterogeneous\_r\_test\_digest}, and
\texttt{general\_policy\_test\_digest}.  Any change to the model code
that would alter a digest is a recovery violation and must either be
reverted or accompanied by a new recovery contract; the rebuild
contract permits no extension whose limit does not recover the
inherited model.

% ---------------------------------------------------------------------
\subsection{Simulation protocol}
\label{sec:methods-sims}

Each simulation under \texttt{Rebuild/sims/} follows a uniform
protocol so that any rebuilt claim can be reproduced byte-for-byte
from the artifact catalogue alone:

\begin{enumerate}
  \item \textbf{Seeded parameter grid.}  Every stochastic component
    of the simulation (random restarts of the policy optimiser,
    multi-start coordinate ascent for the $\Nloc$-D simplex,
    Gauss--Hermite node positioning) is seeded with the run id at the
    head of the script.  The full parameter grid --- ranges of
    $(\valid, \val, \Rsens, \corr, p)$ together with the optimiser
    tolerances --- is recorded both in the simulation's
    \texttt{README.md} and as the \texttt{metadata.grid} object of
    its \texttt{output/results.json}.
  \item \textbf{Recovery test before sweep.}  Every simulation that
    touches an extension axis re-runs the corresponding row of
    Table~\ref{tab:methods-recovery} before its main sweep, and
    refuses to produce output unless the recovery test passes.
  \item \textbf{Canonical JSON output.}  Sweep results are emitted to
    \texttt{output/results.json} via a sort-keys, indent-$2$ JSON
    dump.  A SHA-256 digest is computed over the JSON content
    (excluding non-deterministic wall-clock fields) and recorded in
    the simulation's \texttt{README.md}; for sims with auxiliary
    payloads (e.g.\ source-payload SHAs for derived sweeps), each
    auxiliary digest is recorded alongside.  Re-running the script
    on the same seed and parameter grid must reproduce the digest
    byte-for-byte.
  \item \textbf{Captioned figures.}  Every simulation that backs a
    manuscript claim produces a figure under \texttt{output/} that the
    manuscript includes verbatim, with a caption stating the claim
    the figure supports and the parameter cell the figure was
    generated at.
  \item \textbf{Backlog cross-reference.}  Every simulation's
    \texttt{README.md} states which row of \texttt{CLAIM\_LEDGER.md}
    it backs and which queued sharpening tasks under
    \texttt{REBUILD\_BACKLOG.md} would refine its result (typically
    a finer grid, a wider $\corr$ envelope, or a variant-B
    re-evaluation).
\end{enumerate}

Table~\ref{tab:methods-sims-summary} lists the nine simulations wired
into the manuscript so far.

\begin{table}[tb]
  \centering
  \small
  \begin{tabular}{@{}llp{2.6cm}p{6cm}@{}}
    \toprule
    Sim directory & Backs & Output digest (SHA-256, prefix) & Manuscript cross-reference \\
    \midrule
    \texttt{A1--rho-channel/}
      & A1 (model)
      & \texttt{b692c064\ldots}
      & Section~\ref{sec:model-rho-channel};
        Figures \ref{fig:vda-rho-variantA}, \ref{fig:vda-rho-variantB},
        \ref{fig:cf-vs-rho} \\
    \texttt{C1--cf-distribution/}
      & C1 ($\CF$ distrib.)
      & \texttt{91fc4692\ldots}
      & Section~\ref{sec:results-c1};
        Table \ref{tab:cf-distribution};
        Figures \ref{fig:cf-histogram}, \ref{fig:cf-heatmap},
        \ref{fig:cf-curves} \\
    \texttt{C2--vda-vs-r-vfamily/}
      & C2 ($\rdagger(\val)$)
      & \texttt{09ecef3c\ldots}
      & Section~\ref{sec:results-c2};
        Equation~\eqref{eq:r-dagger};
        Figures \ref{fig:vda-curves-vfamily}, \ref{fig:r-dagger-vs-v} \\
    \texttt{C3--iso-vda-Vv/}
      & C3 (graded band)
      & \texttt{72820559\ldots}
      & Section~\ref{sec:results-c3};
        Table \ref{tab:c3-highV-probe};
        Figures \ref{fig:iso-vda-contours}, \ref{fig:vda-at-high-V} \\
    \texttt{C4--anti-cue-inversion/}
      & C4 (anti-cue)
      & \texttt{6ad651d6\ldots}
      & Section~\ref{sec:results-c4};
        Equation~\eqref{eq:value-weight};
        Table \ref{tab:c4-anticue};
        Figures \ref{fig:c4-alpha-star-map}, \ref{fig:c4-rinv-closed-form} \\
    \texttt{A1--vda-signflip-cellwise/}
      & A1 (cell-wise)
      & \texttt{489c7c25\ldots}
      & Section~\ref{sec:model-upper-bound};
        Table \ref{tab:a1cw-summary};
        Figures \ref{fig:a1cw-sign-heatmap-v5},
        \ref{fig:a1cw-signflip-by-r},
        \ref{fig:a1cw-delta-distribution} \\
    \texttt{A2--heterogeneous-r/}
      & A2 (band)
      & \texttt{22b183f9\ldots}
      & Section~\ref{sec:extensions-a2};
        Table \ref{tab:a2-rb021-summary};
        Figures \ref{fig:a2-vda-peak-band}, \ref{fig:a2-vda-curves-spread} \\
    \texttt{A3--conservation-band/}
      & A3 (band)
      & \texttt{055bf4ec\ldots}
      & Section~\ref{sec:extensions-a3};
        Theorem~\ref{thm:delta-cf-monotone};
        Tables \ref{tab:a3-c1-cf-band}, \ref{tab:a3-c2-peak-band};
        Figures \ref{fig:a3-vda-peak-band}, \ref{fig:a3-delta-cf-distribution} \\
    \texttt{A8--nd-uncued-sweep/}
      & A8 ($\Nloc$-D)
      & \texttt{beb2aa87\ldots}
      & Section~\ref{sec:extensions-a8};
        Table \ref{tab:a8-rb027-summary};
        Figures \ref{fig:a8-simplex-dr}, \ref{fig:a8-anticued-suppression} \\
    \bottomrule
  \end{tabular}
  \caption[Wired simulations]{The nine simulations under
    \texttt{Rebuild/sims/} that back the manuscript's quantitative
    claims.  Each row gives the simulation directory, the claim it
    backs, the leading bytes of the SHA-256 digest of its canonical
    JSON output, and the manuscript cross-references that consume its
    figures and tables.  Full digests are recorded verbatim in each
    simulation's \texttt{README.md} and in
    \texttt{Rebuild/rebuilder\_state.json}.}
  \label{tab:methods-sims-summary}
\end{table}

% ---------------------------------------------------------------------
\subsection{Derivations}
\label{sec:methods-derivations}

Four derivations under \texttt{Rebuild/derivations/} establish in the
rebuild's voice the closed-form results that the body sections cite.
Each derivation is self-contained (setup, exact formulation, full
proof, scope), independently authored from the skeptical reviewer's
attack derivations, and verified by a finite-difference cross-check
or by reduction to a published identity:

\begin{itemize}
  \item \textbf{\texttt{A1--rho-channel.md}} (rb-008,
    Appendix~\ref{sec:appendix-deriv-a1}).  Promotes the
    Booking-1 one-factor reduction of the equicorrelated-Gaussian
    $\PnoFA$ to a derivation in the rebuild's voice; states Slepian
    monotonicity (\cite{Slepian1962}, Tong 1990) at proposition
    strength; gives the two-channel sign decomposition of
    $\partial \VDA / \partial \corr$ that backs the
    Section~\ref{sec:model-upper-bound} retraction of the inherited
    paper's \S5.5 ``upper bound on VDA'' self-characterisation;
    verified against the recovery digest
    \texttt{d3c62215\ldots} of Table~\ref{tab:methods-recovery}.

  \item \textbf{\texttt{C2--non-monotonic-vda-rho.md}} (rb-014,
    Appendix~\ref{sec:appendix-deriv-c2}).  Derives the closed-form
    escape threshold $\rdagger(\val) = K_u/[(\Nloc - 1) K_c]$
    (Equation~\eqref{eq:r-dagger}) by left-derivative of expected
    reward at $\alphacued = 1/\Nloc$, with $\corr > 0$ extension
    via the $\corr$-aware gradient quadratures
    $K_c(\valid, \val; \corr)$ and $K_u(\valid, \val; \corr)$
    (Equations~\eqref{eq:c2-Kc-rho}, \eqref{eq:c2-Ku-rho}) giving
    Proposition~\ref{prop:r-dagger-rho}.  Backs
    Section~\ref{sec:results-c2}; verified numerically at the
    headline cell to floating-point identity against the C2 sim
    digest \texttt{09ecef3c\ldots}.

  \item \textbf{\texttt{C4--anti-cue-inversion.md}} (rb-026,
    Appendix~\ref{sec:appendix-deriv-c2}, downstream).  Derives the
    conditional theorem $\valid \ge 1 / [(\Nloc - 1) \val + 1]$
    (Equation~\eqref{eq:value-weight}) governing whether the cued
    location carries weight at least equal to a uniform allocation,
    and the closed-form inversion threshold $\rstarinv$
    (Equation~\eqref{eq:r-inv}) at which the optimal allocation
    flips to put \emph{less} attention at the cued location under
    counter-predictive cues $\valid < 1/\Nloc$.  Backs
    Section~\ref{sec:results-c4}; verified at the $\Nloc = 4$
    anti-cue cells against the C4 sim digest \texttt{6ad651d6\ldots}.

  \item \textbf{\texttt{A3--power-mean-conservation.md}} (rb-029,
    Appendix~\ref{sec:appendix-deriv-a3}).  Establishes the
    power-mean conservation family as the canonical generalisation of
    A3; gives the boxed Hardy--Littlewood--P\'olya monotonicity in
    the closed KL-divergence form
    $\partial \ln \cost / \partial p = -(1/p^2) \DKL(\mathrm{Bern}(\theta_p) \| \mathrm{Bern}(1/2))$
    (Equation~\eqref{eq:a3d-hlp-kl}); proves
    Proposition~\ref{prop:r-dagger-invariance} ($p$-invariance of
    $\rdagger(\val)$) in full and Corollary~\ref{cor:a3d-c5-invariance}
    (C5 conservation-form-invariance).  Backs
    Section~\ref{sec:extensions-a3} and
    Appendix~\ref{sec:appendix-c5}; verified by finite-difference
    cross-check of Equation~\eqref{eq:a3d-hlp-kl} at seven test
    pairs with LHS--RHS agreement $\le 1.5 \times 10^{-10}$.
\end{itemize}

% ---------------------------------------------------------------------
\subsection{Reproducibility and the rebuild contract}
\label{sec:methods-reproducibility}

The rebuilt manuscript reports a result only when three artifacts
co-exist: a recovery contract that pins the relevant extension to the
inherited model in the appropriate limit; a deterministic simulation
whose output digest is recorded verbatim in
\texttt{Rebuild/rebuilder\_state.json}; and, for every novel
proposition, a derivation under \texttt{Rebuild/derivations/}.  This
\emph{simulate-first, write-second} protocol --- the rebuild's
operating mode, as set out in the mission file --- is what licenses
the manuscript's distributional, graded, and conditional voice
(Section~\ref{sec:limitations-not-claimed}): when a simulation
\emph{weakens} a rebuilt claim, the claim moves down to its supported
strength rather than being suppressed.  No quantitative claim in the
body sections exceeds the strength its row of
\texttt{CLAIM\_LEDGER.md} licenses, and that row in turn points at
the recovery digest, simulation digest, and (where novel) derivation
file that establish the claim's support.

Software dependencies are recorded in \texttt{Rebuild/manuscript/BUILD.md}:
NumPy, SciPy (with a hand-rolled \texttt{erf} fallback for environments
in which SciPy is unavailable), and Matplotlib for the simulation
side; \texttt{pdflatex}, \texttt{bibtex}, and the AMS-LaTeX packages
\texttt{amsmath}, \texttt{amssymb}, \texttt{amsthm}, \texttt{mathtools},
\texttt{booktabs}, and \texttt{hyperref} for the manuscript side.
Every simulation is decomposed into a single \texttt{run.py} that
completes within $\sim 10\text{--}20$ minutes of wall-clock on a
laptop; larger sweeps are decomposed into multiple runs under the
backlog rather than rolled up into a single multi-hour job.

The full reproducibility ledger --- model recovery digests, per-sim
output digests, manuscript PDF byte counts, and the disposition of
every backlog task --- is maintained in
\texttt{Rebuild/rebuilder\_state.json} (atomically rewritten at the
end of every run) and chronicled in \texttt{Rebuild/BUILD\_LOG.md}
(append-only, newest at top).  Cross-claim disposition --- one row per
inherited claim or assumption, the rebuilt strength, and the
backing artifacts --- is maintained in \texttt{Rebuild/CLAIM\_LEDGER.md}.
A reader who wishes to verify any quantitative statement in this
manuscript may locate the claim in \texttt{CLAIM\_LEDGER.md}, follow
the pointer to its backing simulation or derivation, and re-run the
simulation to reproduce the digest byte-for-byte.
